% English mathematical concept atlas: foundational concepts F01--F34.
% Included by the root document; the root defines \concept and math macros.

\concept{F01}{Propositions and quantifiers}{None}
\paragraph{Foundation and definition.}
A proposition is a statement with a truth value after its terms have been
specified. A predicate is a statement that still depends on a variable.
The universal quantifier $\forall$ means ``for every''; the existential
quantifier $\exists$ means ``there is at least one.'' These are logical
definitions, rather than empirical claims about sensor data.
\paragraph{Symbols and meaning.}
Let $\mathcal S$ be a set of permitted input windows, let $x$ denote one
window in that set, let $h:\mathcal S\to\{1,\ldots,K\}$ be a deterministic
classifier, let the integer $K\ge2$ count classes, and let $c$ be one fixed
class. The predicate $P(x)$ is the statement $h(x)=c$. A window can be
represented by $768$ real coordinates, but the logical statement does not
depend on that particular dimension.
\paragraph{Derivation or example.}
The statements
\[
 \forall x\in\mathcal S\;P(x),\qquad
 \exists x\in\mathcal S\;P(x)
\]
have very different strengths. For $\mathcal S=\{u,v\}$, with
$h(u)=c$ and $h(v)\ne c$, the existential statement is true and the
universal statement is false. Negation changes the quantifier:
\[
 \neg(\forall x\in\mathcal S\;P(x))
 \ \Longleftrightarrow\ 
 \exists x\in\mathcal S\;\neg P(x).
\]
Thus a single counterexample defeats a universal claim. Testing a finite
selection from an infinite region does not establish that claim.
\paragraph{Why ECDR needs it.}
A regional certificate asserts that every input admitted by its domain
has the certified decision. Observing that several neighboring windows
share a prediction supplies examples, not the missing universal argument.
\paragraph{Conditions and connections.}
Quantifier order matters: ``for every input there exists a safe radius''
does not specify one radius that works for every input. Universal
statements over the empty set are vacuously true; a useful reuse domain
must also be nonempty. F02 explains the logical role of a certificate.

\concept{F02}{Implication, necessity, sufficiency, and equivalence}{F01}
\paragraph{Foundation and definition.}
An implication $A\Rightarrow B$ says that whenever proposition $A$ holds,
proposition $B$ holds. Then $A$ is sufficient for $B$, and $B$ is necessary
for $A$. Equivalence $A\Leftrightarrow B$ requires both implications.
These relationships describe logical strength; they do not assert that
either proposition is actually true in a particular run.
\paragraph{Symbols and meaning.}
Let $A$ mean ``a sound verifier returns a strictly positive lower bound,''
and let $B$ mean ``the chosen class wins throughout the specified input
region.'' The word sound means that the verifier's claimed inequality is
mathematically valid under its stated model and arithmetic assumptions.
Neither letter is a numerical variable or a dataset label.
\paragraph{Derivation or example.}
For a scalar variable $x\in\R$, take the margin $m(x)=2-x$ on the
interval $\mathcal S=[0,1]$. Its exact minimum is $1$. A valid but loose
lower bound can be $-3$, because $-3\le m(x)$ for every permitted $x$.
Therefore a positive returned bound proves safety, but a nonpositive
returned bound need not prove unsafety:
\[
 \text{positive sound lower bound}\Rightarrow\text{safety},
 \qquad
 \text{no positive bound}\not\Rightarrow\text{unsafety}.
\]
The contrapositive $\neg B\Rightarrow\neg A$ follows from
$A\Rightarrow B$. The converse $B\Rightarrow A$ does not follow.
\paragraph{Why ECDR needs it.}
An unsuccessful certificate attempt should trigger the defined fallback.
It must not be reported as a discovered decision error. Conversely, a
successful attempt is useful only if its soundness conditions hold.
\paragraph{Conditions and connections.}
A conditional theorem proves its conclusion from hypotheses. An
unconditional claim has no remaining hypotheses within its declared
framework; it still uses the framework's axioms and definitions. Saying
``unconditional guarantee'' cannot remove model identity, input-domain,
or numerical assumptions. F28 gives a concrete conditional theorem.

\concept{F03}{Sets, membership, and inclusion}{F01}
\paragraph{Foundation and definition.}
A set specifies which objects belong to a collection. Membership
$x\in\mathcal S$ is a yes-or-no relation. Inclusion
$\mathcal A\subseteq\mathcal B$ means every member of $\mathcal A$ is
also a member of $\mathcal B$. A set is different from its size, its
bounding box, and a list of sampled members.
\paragraph{Symbols and meaning.}
Here $x$ is an input vector; $\mathcal S$, $\mathcal A$, and
$\mathcal B$ are sets of input vectors in the same coordinate system.
The empty set $\varnothing$ has no members. The union
$\mathcal A\cup\mathcal B$ admits a member of either set; the
intersection $\mathcal A\cap\mathcal B$ admits only shared members.
The difference $\mathcal B\setminus\mathcal A$ contains members of
$\mathcal B$ outside $\mathcal A$.
\paragraph{Derivation or example.}
Let $\mathcal A=[-1,1]$ and $\mathcal B=[-2,2]$, subsets of $\R$.
Then $0.5\in\mathcal A\subseteq\mathcal B$, whereas
$1.5\in\mathcal B\setminus\mathcal A$. If a classifier $h$ satisfies
$h(x)=c$ for all $x\in\mathcal B$, it also satisfies this identity on
$\mathcal A$. Proof: choose any member $x$ of $\mathcal A$; inclusion
puts the same $x$ in $\mathcal B$, where the identity is known.
The reverse transfer is invalid without more information.
\paragraph{Why ECDR needs it.}
Membership is the online gate that connects a particular query to a
previously proved regional statement. A conservative inner gate can
reduce reuse coverage while preserving safety. Enlarging the gate beyond
the certified set can invalidate safety.
\paragraph{Conditions and connections.}
An inequality can define a set without proving that observed data usually
falls inside it. All comparisons require identical normalization and
indexing. F24 constructs boxes; F34 explains why individually chosen
coordinate quantiles do not automatically determine joint coverage.

\concept{F04}{Functions, domains, images, and preimages}{F03}
\paragraph{Foundation and definition.}
A function assigns exactly one output to every input in its domain.
Writing $f:\mathcal X\to\mathcal Y$ names the domain $\mathcal X$
and codomain $\mathcal Y$. The image is the set of outputs actually
attained. A preimage collects inputs whose outputs belong to a specified
output set; it does not require an inverse function.
\paragraph{Symbols and meaning.}
The variable $x\in\mathcal X$ is one input and $f(x)\in\mathcal Y$
is its output. For subsets $\mathcal A\subseteq\mathcal X$ and
$\mathcal B\subseteq\mathcal Y$, define
\[
 f(\mathcal A)=\{f(x):x\in\mathcal A\},\qquad
 f^{-1}(\mathcal B)=\{x\in\mathcal X:f(x)\in\mathcal B\}.
\]
The superscript $-1$ here denotes a preimage operation, not division.
\paragraph{Derivation or example.}
For $f:\R\to\R$ given by $f(x)=x^2$, the codomain is $\R$ but
the image is $[0,\infty)$. The preimage of $[1,4]$ is
$[-2,-1]\cup[1,2]$. No single-valued inverse on the full image can
recover whether the original input was positive or negative.
For a classifier $h:\R^{d}\to\{1,\ldots,K\}$, the decision region
of class $c$ is $h^{-1}(\{c\})$. Here integers $d\ge1$ and $K\ge2$
count input coordinates and classes.
\paragraph{Why ECDR needs it.}
Certification seeks a practical input set contained in one decision
region. Feature compression can map different inputs to the same feature
vector, so invertibility is unnecessary for classification or reuse.
\paragraph{Conditions and connections.}
The actual reference must be deterministic under its inference state;
random layers require a separately stated treatment. F17 explains how
preprocessing, a CNN, and a decision rule form a composed function.

\concept{F05}{Real-number field axioms and algebra}{F02}
\paragraph{Foundation and definition.}
The field axioms are assumptions about addition and multiplication on
$\R$. Each operation is closed. Both are associative and commutative.
There are identities $0$ and $1$ with $0\ne1$. Every element has an
additive inverse, every nonzero element has a multiplicative inverse, and
multiplication distributes over addition. These axioms are not proved
by collecting sensor measurements.
\paragraph{Symbols and meaning.}
For arbitrary real numbers $a,b,c$, associativity includes
$(a+b)+c=a+(b+c)$; distributivity is $a(b+c)=ab+ac$. The symbols
$-a$ and $a^{-1}$ mean the respective inverses, with $a^{-1}$ defined
only for $a\ne0$. Define $a-b=a+(-b)$ and $a/b=a b^{-1}$ for $b\ne0$.
\paragraph{Derivation or example.}
Cancellation follows from the axioms rather than being another empirical
observation. If $a+c=b+c$, adding $-c$ to both sides and reassociating
gives $a=b$. Also,
\[
 a\cdot0=a(0+0)=a\cdot0+a\cdot0
 \quad\Rightarrow\quad a\cdot0=0.
\]
This proves why division by zero cannot be defined as a multiplicative
inverse: an inverse would imply $1=0$. For a logit pair $u,v\in\R$,
the margin $u-v$ is simply addition with the additive inverse of $v$.
\paragraph{Why ECDR needs it.}
Rearranging margin expressions, canceling shared bias terms, and
distributing sums over a linear head all rely on these algebraic laws.
Stating the foundation makes clear which manipulations are legitimate.
\paragraph{Conditions and connections.}
Floating-point arithmetic only approximates this field: addition can
fail associativity after rounding. Real-number proofs therefore need an
additional error argument before becoming machine-level certificates.
Order and completeness require further axioms, developed in F06--F07.

\concept{F06}{Ordered-field axioms and inequalities}{F05}
\paragraph{Foundation and definition.}
An ordered field adds a total order compatible with field operations.
For any two real numbers exactly one of $a<b$, $a=b$, or $b<a$ holds.
Order is transitive; adding the same number preserves order; and the
product of positive numbers is positive. These are axioms of the real
number system, not assumptions about the classifier.
\paragraph{Symbols and meaning.}
The symbols $a,b,c\in\R$ denote arbitrary scalars. Write $a\le b$
when $a<b$ or $a=b$. A positive scalar satisfies $c>0$; a negative scalar
satisfies $c<0$. The notation $[a,b]$ later denotes all scalars between
the endpoints when $a\le b$.
\paragraph{Derivation or example.}
If $a\le b$ and $c>0$, then $ac\le bc$: from $b-a\ge0$, product
compatibility gives $c(b-a)\ge0$, and distributivity finishes the proof.
For $c<0$, multiplying reverses the direction. For example,
\[
 -2\le3\quad\Rightarrow\quad 4\le-6
 \quad\text{is false},\qquad -2\le3\Rightarrow4\ge-6.
\]
If a margin satisfies $m\ge\ell$ and a proved bound gives $\ell>0$,
transitivity yields $m>0$. If only $\ell\le0$ is known, neither the
sign of $m$ nor a decision error follows.
\paragraph{Why ECDR needs it.}
The direction of a bound determines its usefulness. A lower margin bound
supports a safety claim; an upper bound with the same numerical value
generally does not. Negative coefficients also determine which interval
endpoint minimizes an affine expression.
\paragraph{Conditions and connections.}
Order is not a floating-point tolerance policy. Replacing $m>0$ by
$m\ge-10^{-6}$ changes the mathematical acceptance condition and needs
its own justification. F25 distinguishes bounding values from attained
minima; F28 uses strict positivity.

\concept{F07}{Completeness and the least-upper-bound axiom}{F03, F06}
\paragraph{Foundation and definition.}
Completeness distinguishes the real numbers from an arbitrary ordered
field. Its least-upper-bound formulation states: every nonempty subset
of $\R$ that is bounded above has a least upper bound in $\R$. This is
an axiom characterizing the real-number setting. It does not say that
every set contains its upper bound.
\paragraph{Symbols and meaning.}
Let $\mathcal A\subseteq\R$ be nonempty. A number $u\in\R$ is an
upper bound when $a\le u$ for every $a\in\mathcal A$. Its supremum
$s=\sup\mathcal A$ is an upper bound satisfying $s\le u$ for every
upper bound $u$. Bounded above means that at least one such finite $u$
exists.
\paragraph{Derivation or example.}
For $\mathcal A=(0,1)$, the supremum is $1$, although
$1\notin\mathcal A$. Every number below $1$ fails to be an upper bound,
because there is a larger member still below $1$.
Completeness also gives an infimum for a nonempty set bounded below:
\[
 \inf\mathcal A=-\sup(-\mathcal A),\qquad
 -\mathcal A=\{-a:a\in\mathcal A\}.
\]
The ordered rational field lacks this completeness: the positive
rational numbers with square below $2$ have no rational least upper
bound. A real square root supplies the missing bound.
\paragraph{Why ECDR needs it.}
Worst-case optimization uses infima over infinitely many possible inputs.
Completeness supports the real-valued limits and extrema theory behind
certification, even when a practical verifier returns only an enclosure.
\paragraph{Conditions and connections.}
Completeness does not supply an algorithm or guarantee finite computation.
An unbounded margin set may have extended-real infimum $-\infty$ instead
of a finite real bound. F25 explains infimum versus minimum, and
F27--F28 add hypotheses that make an extremum attainable.

\concept{F08}{Integer indexing, floor, and ceiling}{F05, F06}
\paragraph{Foundation and definition.}
An index identifies a discrete position. For a real scalar $u$,
$\lfloor u\rfloor$ is the greatest integer not exceeding $u$, and
$\lceil u\rceil$ is the least integer not smaller than $u$. These
definitions convert continuous counts or thresholds into valid discrete
positions; choosing zero-based indexing is a convention, not a theorem.
\paragraph{Symbols and meaning.}
Let the integers $T\ge1$, $W\ge1$, and $s\ge1$ count recording samples,
window samples, and stride samples. Let $i$ be a zero-based window index.
Window $i$ starts at sample $is$ and contains indices
$is,is+1,\ldots,is+W-1$. A sampling frequency $f_s>0$, measured in
samples per second, converts index differences to seconds.
\paragraph{Derivation or example.}
A full window exists exactly when $is+W\le T$. Hence its count is
\[
 N=\max\!\left(0,1+\left\lfloor\frac{T-W}{s}\right\rfloor\right).
\]
For $T=1000$, $W=128$, and $s=64$, there are $N=14$ windows, with
final start $13\cdot64=832$ and final included sample $959$.
The remaining $40$ samples are outside this grid's windows. If a class
must occupy at least fraction $\tau=0.6$ of a $128$-sample window,
the required integer count is $\lceil\tau W\rceil=77$.
\paragraph{Why ECDR needs it.}
The overlap structure, valid CNN output length, and available transition
examples all depend on exact indexing. A one-position convention change
can change labels or invalidate an equality used to reuse activations.
\paragraph{Conditions and connections.}
An inclusive annotation interval $[a,b]$ contains $b-a+1$ integer
positions; a half-open slice $[a,b)$ contains $b-a$. Neither convention
should be inferred from notation alone. The index extent $W/f_s$ and
first-to-last timestamp gap $(W-1)/f_s$ are distinct time quantities.

\concept{F09}{Finite sums and reindexing}{F05, F08}
\paragraph{Foundation and definition.}
A finite sum repeatedly adds terms over a declared finite index set.
Associativity and commutativity of real addition justify regrouping and
reordering. Reindexing is valid when the new indices describe exactly the
same terms once each, or when explicitly stated terms are added or removed.
\paragraph{Symbols and meaning.}
Let $n\ge1$ be an integer, let $a_0,\ldots,a_{n-1}\in\R$ be scalars,
and let $i$ be a summation index. Then $\sum_{i=0}^{n-1}a_i$ adds those
$n$ terms. For vector terms the same notation means coordinatewise
addition. A sum over an empty index set is defined as the additive
identity, either zero or the zero vector as appropriate.
\paragraph{Derivation or example.}
Let $H_0,\ldots,H_{n+s-1}\in\R^d$ be feature vectors, with integers
$1\le s\le n$ and $d\ge1$. Reindexing gives
\[
 \sum_{i=0}^{n-1}H_{i+s}-\sum_{i=0}^{n-1}H_i
 =\sum_{i=n}^{n+s-1}H_i-\sum_{i=0}^{s-1}H_i.
\]
Both left-hand sums contain the positions $s,\ldots,n-1$, whose terms
cancel. For $n=4,s=2$, this is
$(H_2+H_3+H_4+H_5)-(H_0+H_1+H_2+H_3)
=H_4+H_5-H_0-H_1$. Dividing by $n$ yields a difference of averages.
\paragraph{Why ECDR needs it.}
This cancellation explains why a global-average-pooled feature can change
through entering and leaving feature positions when alignment is exact.
It is an algebraic consequence, not an empirical statement that all
nearby windows have small differences.
\paragraph{Conditions and connections.}
The shared symbols must represent equal features, established by a
separate receptive-field and shift argument. Equal input-window lengths
alone do not establish that equality. In floating-point arithmetic,
different summation orders can produce different rounded outputs.

\concept{F10}{Vector spaces and their axioms}{F05}
\paragraph{Foundation and definition.}
A real vector space is a set $V$ equipped with vector addition and real
scalar multiplication. Addition is associative and commutative, has a
zero vector, and supplies an additive inverse for every vector. Scalar
multiplication satisfies identity, compatibility, and both distributive
laws. Closure requires the resulting objects to remain in $V$.
\paragraph{Symbols and meaning.}
Let $u,v,w\in V$ be vectors and $a,b\in\R$ be scalars. The remaining
laws are $1v=v$, $(ab)v=a(bv)$,
$a(u+v)=au+av$, and $(a+b)v=av+bv$. In the concrete space $\R^d$,
the positive integer $d$ counts coordinates, and these operations act
coordinatewise. The zero vector is written $0_V$ when confusion with
the real scalar zero is possible.
\paragraph{Derivation or example.}
For $u=(1,2)$ and $v=(3,-1)$ in $\R^2$,
$2u-v=(-1,5)$. Addition is legitimate because corresponding entries use
the same coordinate meanings. A set of physically realizable sensor
windows need not itself be a vector space: adding two realizable windows
or multiplying one by a negative number may leave the physical set.
The ambient coordinate space can still be a vector space containing it.
\paragraph{Why ECDR needs it.}
Differences between windows, linear heads, and perturbation directions
are vectors in an ambient space. This makes linear algebra available
without assuming that every resulting vector is a plausible activity.
\paragraph{Conditions and connections.}
Coordinate scaling must be fixed when physical channels have different
units. An acceleration coordinate in $g$ and an angular-velocity coordinate
in rad/s are distinct types of measurements. Normalized coordinates
support a declared dimensionless geometry. F11--F14 develop combinations,
span, independence, and coordinates from these axioms.

\concept{F11}{Linear combinations}{F09, F10}
\paragraph{Foundation and definition.}
A linear combination forms a new vector by scaling finitely many vectors
and adding the results. It is a definition built from the vector-space
operations. Negative coefficients and coefficients whose sum differs
from one are allowed; convex and affine combinations impose additional
constraints and are different concepts.
\paragraph{Symbols and meaning.}
Let the integers $d\ge1$ and $r\ge1$ count coordinates and supplied
directions. Let $v_1,\ldots,v_r\in\R^d$ be direction vectors and let
$\alpha_1,\ldots,\alpha_r\in\R$ be scalar coefficients. Their linear
combination is
\[
 v=\sum_{j=1}^{r}\alpha_j v_j.
\]
Here $j$ is an index, and $v\in\R^d$ is the resulting vector.
If the directions are dimensionless and the coefficients are
dimensionless, so is the combination.
\paragraph{Derivation or example.}
Take $v_1=(1,1)$ and $v_2=(1,-1)$. The coefficients
$\alpha_1=2$ and $\alpha_2=1$ give $v=(3,1)$.
Collecting the directions as columns of
$B=\begin{pmatrix}1&1\\1&-1\end{pmatrix}$ and coefficients as
$\alpha=(2,1)^\top$ expresses the same operation as $v=B\alpha$.
Restricting coefficients to $|\alpha_1|\le2$ and $|\alpha_2|\le1$
selects a bounded parallelogram from all possible combinations.
\paragraph{Why ECDR needs it.}
A structured perturbation model can represent correlated changes using
several directions instead of an independent allowance at every input
coordinate. The coefficient bounds specify how much motion along each
direction is admitted.
\paragraph{Conditions and connections.}
Representability does not establish that a direction is physically
meaningful, statistically likely, or safe for the classifier. Those are
separate modeling, empirical, and verification questions. F12 allows all
coefficients to define a span; a bounded coefficient box gives only a
subset of that span, which must not be confused with the whole subspace.

\concept{F12}{Linear span: all combinations and the smallest subspace}{F10, F11}
\paragraph{Foundation and definition.}
The linear span of vectors is the set of every vector obtainable as a
linear combination. For directions $v_1,\ldots,v_r\in\R^d$, define
\[
 \operatorname{span}\{v_1,\ldots,v_r\}
 =\left\{\sum_{j=1}^{r}\alpha_jv_j:
                   \alpha_1,\ldots,\alpha_r\in\R\right\}.
\]
The positive integers $r$ and $d$ count supplied directions and ambient
coordinates. All coefficients $\alpha_j$ are unrestricted real numbers.
This linear-algebra term differs from a temporal annotation span.
\paragraph{Symbols and meaning.}
A subspace is a subset containing the zero vector and closed under
addition and real scaling. Here $j$ indexes directions; a member of the
span is a possible direction combination. The word ``linear'' implies
that the resulting set passes through the origin.
\paragraph{Derivation or example.}
For $v_1=(1,2)$, its span is the line
$\{(a,2a):a\in\R\}$. Adding $v_2=(0,1)$ produces the entire plane:
an arbitrary $(x,y)\in\R^2$ equals $xv_1+(y-2x)v_2$.
The span is the \emph{smallest} subspace containing the supplied vectors.
Proof: zero coefficients give zero; adding combinations adds their
coefficients; scaling scales the coefficients, so the span is a subspace.
Any subspace containing the supplied vectors must contain every scaled
vector and every finite sum of them, hence contains the span.
\paragraph{Why ECDR needs it.}
If a matrix $B$ contains learned perturbation directions as columns,
$B\alpha$ stays in their span. A residual term is needed to represent
changes outside that subspace; discarding it without evidence can exclude
real queries from the modeled region.
\paragraph{Conditions and connections.}
Dependent directions do not increase span dimension. Bounding
coefficients changes an unbounded subspace into a restricted region.
Translating by an anchor creates an affine set. F13--F16 explain these
distinctions; F19 explains orthogonal residuals.

\concept{F13}{Linear independence and redundant directions}{F11, F12}
\paragraph{Foundation and definition.}
Vectors are linearly independent when the only linear combination equal
to zero uses zero for every coefficient. Otherwise they are dependent.
Independence is an algebraic property. It does not mean probabilistic
independence, uncorrelated measurements, or distinct class labels.
\paragraph{Symbols and meaning.}
For vectors $v_1,\ldots,v_r\in\R^d$, with integers $r,d\ge1$, and
coefficients $\alpha_1,\ldots,\alpha_r\in\R$, independence means
\[
 \sum_{j=1}^{r}\alpha_jv_j=0
 \quad\Rightarrow\quad
 \alpha_1=\cdots=\alpha_r=0.
\]
Here $0$ on the left is the $d$-dimensional zero vector, and $j$ indexes
directions. A nonzero coefficient vector producing zero demonstrates
redundancy among the directions.
\paragraph{Derivation or example.}
The vectors $v_1=(1,2)$ and $v_2=(2,4)$ are dependent because
$2v_1-v_2=0$. Their span remains one line despite two supplied columns.
In contrast, $u_1=(1,0)$ and $u_2=(1,1)$ are independent: if
$a u_1+b u_2=(a+b,b)=(0,0)$, then $b=0$ and $a=0$.
Independent directions also give unique coefficients for any vector in
their span. If two coefficient lists represent the same vector,
subtracting them gives a zero combination; independence forces each
coefficient difference to vanish.
\paragraph{Why ECDR needs it.}
An independent basis prevents ambiguity when coefficients describe a
structured perturbation. With redundant columns, the same perturbation
can have several coefficient vectors, so testing one arbitrary encoding
against coefficient bounds may not settle set membership.
\paragraph{Conditions and connections.}
Near-dependence is a numerical conditioning issue even when exact
independence holds. A tolerance-based rank estimate should therefore be
reported as a numerical decision. No more than $d$ vectors can be
independent in $\R^d$. F14 uses independence together with spanning to
define a basis.

\concept{F14}{Bases, coordinates, and dimension}{F12, F13}
\paragraph{Foundation and definition.}
A basis of a vector space is a linearly independent set that spans the
space. Every vector then has a unique linear-combination representation.
For a finite-dimensional space, its dimension is the number of elements
in any basis; equality of basis sizes is a theorem, not the definition
of independence.
\paragraph{Symbols and meaning.}
Let $V\subseteq\R^d$ be a subspace of dimension $r$, where integers
$1\le r\le d$. Let $b_1,\ldots,b_r\in\R^d$ be a basis and
$B\in\R^{d\times r}$ the matrix with these vectors as columns.
The coordinate vector $\alpha\in\R^r$ of $v\in V$ satisfies
$v=B\alpha$. The ambient dimension $d$ and subspace dimension $r$ have
different meanings.
\paragraph{Derivation or example.}
Take $b_1=(1,1)$ and $b_2=(1,-1)$ as a basis of $\R^2$.
For $v=(3,1)$, solving
$\alpha_1+\alpha_2=3$ and $\alpha_1-\alpha_2=1$ gives the unique
coordinates $\alpha=(2,1)^\top$. Changing the basis changes these
coordinates without changing $v$. For example, standard-coordinate
axes represent the same vector by $(3,1)^\top$.
\paragraph{Why ECDR needs it.}
A low-dimensional model describes perturbations using $r$ coefficients
instead of $d$ ambient coordinates. This can express measured structure,
but dimensionality alone does not establish classifier invariance or
an efficient online membership algorithm.
\paragraph{Conditions and connections.}
A truncated direction matrix spans a subspace rather than the whole
input space. A residual must account for omitted directions when they
are allowed. For orthonormal columns, coordinates are $B^\top v$;
for a general independent basis they are
$(B^\top B)^{-1}B^\top v$. Using the first formula without
orthonormality generally gives incorrect coordinates.

\concept{F15}{Matrices as linear maps}{F09, F10, F14}
\paragraph{Foundation and definition.}
A map is linear if it preserves vector addition and scalar
multiplication. A real matrix represents a linear map after input and
output bases are chosen. Conversely, each linear map between
finite-dimensional coordinate spaces has a matrix representation.
\paragraph{Symbols and meaning.}
Let $A\in\R^{m\times d}$ be a matrix, $x\in\R^d$ an input, and
$y=Ax\in\R^m$ an output. Positive integers $d$ and $m$ count input
and output coordinates. Entry $A_{ji}$ multiplies input coordinate
$x_i$ when forming output coordinate $y_j$:
\[
 y_j=\sum_{i=1}^{d}A_{ji}x_i,
 \qquad j=1,\ldots,m.
\]
The transpose $A^\top\in\R^{d\times m}$ exchanges row and column
indices; it is not generally an inverse.
\paragraph{Derivation or example.}
For
$A=\begin{pmatrix}2&-1\\0&3\end{pmatrix}$ and $x=(1,2)^\top$,
the output is $(0,6)^\top$. Distributivity proves
$A(x+z)=Ax+Az$ for any $z\in\R^d$: each row distributes over the
corresponding scalar sum. Similarly, $A(a x)=a(Ax)$ for any scalar
$a\in\R$. A map containing a fixed nonzero bias generally fails
linearity because its output at zero is not zero.
\paragraph{Why ECDR needs it.}
A convolution with fixed weights and no bias can be represented as a
sparse matrix. A final linear head is a small dense matrix. This shared
language allows coordinate differences and norm bounds to propagate
through both kinds of layers.
\paragraph{Conditions and connections.}
Shape compatibility is necessary for composition and multiplication.
Different flattening orders require corresponding matrix-column orders.
Physical units must also match: a coefficient transforms input units
into output units. F16 adds biases; F21 measures worst-case amplification.

\concept{F16}{Affine maps and translated structure}{F11, F15}
\paragraph{Foundation and definition.}
An affine map has the form of a linear map plus a constant vector.
It preserves affine combinations, whose scalar coefficients sum to one.
It need not preserve the origin. A translated subspace is an affine set,
which should not be called a linear subspace unless it contains zero.
\paragraph{Symbols and meaning.}
Let $A\in\R^{m\times d}$ be fixed weights,
$b\in\R^m$ a fixed bias, and $x\in\R^d$ an input. Integers
$m,d\ge1$ count output and input coordinates. Define
$f(x)=Ax+b$. For an anchor $x_0\in\R^d$ and perturbation
$\delta\in\R^d$, the new input is $x_0+\delta$.
\paragraph{Derivation or example.}
The difference identity is exact over real numbers:
\[
 f(x_0+\delta)-f(x_0)
 =A(x_0+\delta)+b-(Ax_0+b)=A\delta.
\]
Thus a fixed bias affects absolute outputs but cancels from their
differences. In one dimension, $f(x)=2x+3$ gives
$f(4)-f(1)=11-5=6=2(4-1)$.
If $u,v\in\R^d$ and $\lambda\in\R$, then
$f(\lambda u+(1-\lambda)v)=\lambda f(u)+(1-\lambda)f(v)$,
because the two coefficients of $b$ sum to one.
\paragraph{Why ECDR needs it.}
Class logits and fixed normalization are affine operations. Their
differences can often be analyzed more tightly than separate absolute
output intervals. Translating a perturbation domain by an anchor also
explains precisely where a local certificate applies.
\paragraph{Conditions and connections.}
The same $A$ and $b$ must be used at both inputs. Updating running
normalization statistics or model weights breaks that cancellation.
ReLU layers make a network piecewise affine; one affine formula applies
only when its activation conditions are established.

\concept{F17}{Composition and the complete reference function}{F04, F15, F16}
\paragraph{Foundation and definition.}
Composition applies one function after another. If the first function's
outputs lie within the second function's domain, the composition is
well defined. Writing a pipeline as a composition exposes which changes
alter the mathematical reference being certified.
\paragraph{Symbols and meaning.}
Let $P:\R^{d}\to\R^{d}$ be preprocessing,
$G:\R^{d}\to\R^{r}$ the feature extractor,
$L:\R^{r}\to\R^{K}$ the logit head, and
$D:\R^{K}\to\{1,\ldots,K\}$ the decision rule. The positive
integers $d,r,K$ count input coordinates, features, and classes.
Define $h=D\circ L\circ G\circ P$, so
$h(x)=D(L(G(P(x))))$ for input $x\in\R^d$.
\paragraph{Derivation or example.}
In a scalar example, let $P(x)=(x-2)/4$ and $G(u)=\max(0,u)$.
Then $(G\circ P)(6)=1$, while $(G\circ P)(0)=0$.
Reversing the order gives $(P\circ G)(0)=-0.5$, which differs from
zero. Composition is therefore generally noncommutative. It is
associative when all involved domains fit: grouping stages does not
change the order of evaluation.
\paragraph{Why ECDR needs it.}
A certificate for logits under one channel ordering or normalization is
not automatically a certificate for a changed pipeline. The reference
includes preprocessing, inference state, weights, and deterministic tie
handling, even if only the CNN weights are stored in a checkpoint.
\paragraph{Conditions and connections.}
Every stage must receive the shape, units, and coordinate conventions it
expects. A certificate built in normalized coordinates needs the query
mapped by the same fixed $P$. F21 shows how amplification bounds
multiply through composition; nonlinear stages require their own bounds.

\concept{F18}{Inner products and dot products}{F09, F10}
\paragraph{Foundation and definition.}
A real inner product is a symmetric bilinear operation producing a
scalar and satisfying positive definiteness: a vector's product with
itself is nonnegative and vanishes only at zero. The standard dot
product on $\R^d$ is one particular inner product.
\paragraph{Symbols and meaning.}
For vectors $u,v\in\R^d$, where integer $d\ge1$ counts coordinates,
write
\[
 \langle u,v\rangle=u^\top v=\sum_{i=1}^{d}u_i v_i.
\]
The index $i$ selects a coordinate. Symmetry means
$\langle u,v\rangle=\langle v,u\rangle$; bilinearity means linearity
in either argument when the other is fixed. A weighted product
$u^\top Mv$ is an inner product when the matrix
$M\in\R^{d\times d}$ is symmetric positive definite.
\paragraph{Derivation or example.}
With $u=(1,2)$ and $v=(3,-1)$, their dot product is $3-2=1$.
Also $u^\top u=5$. For a head row difference $q=(2,-1)$ and feature
change $\delta=(0.1,0.3)$, the scalar change is
$q^\top\delta=0.2-0.3=-0.1$. Opposing coordinate contributions can
partly cancel; their separate absolute magnitudes would sum to $0.5$.
\paragraph{Why ECDR needs it.}
A class-pair margin is a directional measurement of a feature vector.
Keeping that direction can preserve cancellations that separate logit
bounds lose. Inner products also define orthogonality, projections, and
the usual geometry of learned directions.
\paragraph{Conditions and connections.}
The standard dot product implicitly chooses a coordinate scale. Adding
products of physically different units needs a declared normalization
or weighting interpretation. F19 develops orthogonality; F22 derives the
dual-norm inequality that bounds a directional change by its magnitude.

\concept{F19}{Orthogonality and residual decomposition}{F18}
\paragraph{Foundation and definition.}
Two vectors are orthogonal when their inner product is zero. This
algebraic relation means a right angle for nonzero vectors under the
chosen geometry. Orthonormal vectors are pairwise orthogonal and each
has unit Euclidean length. Neither condition is equivalent to
probabilistic independence.
\paragraph{Symbols and meaning.}
Let $U\in\R^{d\times r}$ have orthonormal columns, with integers
$1\le r\le d$. Thus $U^\top U=I_r$, where $I_r$ is the
$r\times r$ identity matrix. For a vector $\delta\in\R^d$, define
its coefficient vector $\alpha=U^\top\delta\in\R^r$ and residual
$e=\delta-U\alpha\in\R^d$.
\paragraph{Derivation or example.}
Multiplying the residual by $U^\top$ gives
\[
 U^\top e=U^\top\delta-U^\top U\alpha
 =\alpha-\alpha=0.
\]
The residual is therefore orthogonal to every supplied direction, and
$\delta=U\alpha+e$. For $U=(1,0)^\top$ and $\delta=(3,4)^\top$,
the projected part is $(3,0)^\top$ and the residual is $(0,4)^\top$.
Expanding the dot product gives Pythagoras:
$\norm{\delta}_2^2=\norm{U\alpha}_2^2+\norm{e}_2^2=9+16=25$,
because the cross term vanishes.
\paragraph{Why ECDR needs it.}
A direction-based domain can distinguish explained motion from
unexplained residual motion. Bounding the residual prevents a
low-dimensional approximation from silently asserting that all real
changes lie exactly in the learned subspace.
\paragraph{Conditions and connections.}
Column scaling generally destroys orthonormality, so $U^\top\delta$
cannot then be reused as a coefficient formula. Orthogonality is relative
to an inner product, which must match normalization and membership
calculations. C11 explains projection; C12 explains how PCA estimates
directions from finite data.

\concept{F20}{Norms and the meanings of vector size}{F10, F18}
\paragraph{Foundation and definition.}
A norm assigns a nonnegative size to each vector. Its axioms are
definiteness, absolute homogeneity, and the triangle inequality:
$\norm{v}=0$ exactly at $v=0$,
$\norm{a v}=|a|\norm{v}$, and
$\norm{u+v}\le\norm{u}+\norm{v}$.
Here $u,v$ are vectors in one real vector space and $a\in\R$ is a scalar.
\paragraph{Symbols and meaning.}
For $v\in\R^d$, with integer $d\ge1$ and coordinate index
$i\in\{1,\ldots,d\}$, common choices are
\[
 \norm{v}_1=\sum_{i=1}^d|v_i|,\qquad
 \norm{v}_2=\sqrt{\sum_{i=1}^d v_i^2},\qquad
 \norm{v}_\infty=\max_{1\le i\le d}|v_i|.
\]
The subscripts specify different notions of size; $\infty$ names the
maximum-coordinate norm, not an infinite vector size.
\paragraph{Derivation or example.}
For $v=(3,-4)$, these sizes are $7$, $5$, and $4$. More generally,
\[
 \norm{v}_\infty\le\norm{v}_2\le\norm{v}_1,
 \qquad \norm{v}_2\le\sqrt d\,\norm{v}_\infty.
\]
For example, the last inequality follows by bounding each of the $d$
squared coordinates by $\norm{v}_\infty^2$ and taking square roots.
Thus changing the norm while keeping the same numerical radius changes
the admitted perturbations.
\paragraph{Why ECDR needs it.}
Norms define perturbation budgets, Lipschitz changes, and residual size.
A maximum-coordinate budget constrains every coordinate, whereas an
Euclidean budget constrains total squared change.
\paragraph{Conditions and connections.}
Norm values depend on coordinate scaling. Raw acceleration and angular
velocity should not be treated as physically commensurate just because
their numerical entries share an array. F22 pairs norms with their duals;
F23 turns a norm into a distance and a ball.

\concept{F21}{Induced operator norms and composition bounds}{F15, F20}
\paragraph{Foundation and definition.}
An induced operator norm measures the largest amplification of vector
size by a linear map. For compatible input and output norms, define
\[
 \norm{A}_{p\to q}
 =\sup_{v\ne0}\frac{\norm{Av}_q}{\norm{v}_p}.
\]
Here $A\in\R^{m\times d}$, $v\in\R^d$, integers $m,d\ge1$
count coordinates, and subscripts $p,q$ name the chosen norms.
\paragraph{Symbols and meaning.}
The supremum is the least upper bound over all nonzero input directions.
The notation $\norm{A}_2$ usually means $\norm{A}_{2\to2}$, not
the square root of the sum of every squared matrix entry. The latter is
the Frobenius norm and is a different quantity.
\paragraph{Derivation or example.}
The definition gives $\norm{Av}_q\le\norm{A}_{p\to q}\norm{v}_p$.
For another matrix $B\in\R^{d\times r}$, with integer $r\ge1$ and
input norm indexed by $t$, applying this inequality twice yields
\[
 \norm{AB}_{t\to q}
 \le\norm{A}_{p\to q}\norm{B}_{t\to p}.
\]
This is submultiplicativity. For
$A=\operatorname{diag}(2,-3)$, the Euclidean induced norm is $3$:
each squared output is bounded by $9$ times the squared input, and
input $(0,1)$ attains equality. Under maximum-coordinate norms,
$\norm{A}_{\infty\to\infty}=\max_j\sum_i|A_{ji}|$.
\paragraph{Why ECDR needs it.}
Layerwise amplification bounds can be multiplied to obtain a network
bound. This is simple and sound when every layer constant is valid,
although successive worst-case directions may not be jointly achievable.
\paragraph{Conditions and connections.}
The intermediate norms and dimensions must match. A computed spectral
estimate can underestimate the true norm unless certified with an error
bound. Biases cancel from affine differences; nonlinear stages need
separate Lipschitz arguments.

\concept{F22}{H\"older's inequality and dual norms}{F18, F20}
\paragraph{Foundation and definition.}
The dual norm measures the largest dot product a direction can produce
against a unit-norm vector. For a norm $\norm{\cdot}$ on $\R^d$,
define $\norm{q}_*=\sup_{\norm{v}\le1}|q^\top v|$.
Here integer $d\ge1$ counts coordinates, while $q,v\in\R^d$ are
a direction and a test vector. The star names the dual norm.
\paragraph{Symbols and meaning.}
For $L_p$ norms with $1\le p\le\infty$, the dual is $L_{p^*}$,
where $1/p+1/p^*=1$ and $1/\infty=0$. Thus $L_1$ pairs with
$L_\infty$, $L_2$ pairs with itself, and $L_\infty$ pairs with $L_1$.
The vector $\delta\in\R^d$ below denotes a perturbation.
\paragraph{Derivation or example.}
If $\delta\ne0$, then $v=\delta/\norm{\delta}$ is a unit vector.
Substitution into the dual-norm definition gives
\[
 |q^\top\delta|\le\norm{q}_*\norm{\delta}.
\]
At $\delta=0$ the result follows directly. This is H\"older's
inequality in dual-norm form. If $q=(2,-1)$ and
$\norm{\delta}_\infty\le0.1$, the largest absolute directional
change is $(|2|+|-1|)0.1=0.3$, attained at
$\delta=(0.1,-0.1)$ or its negative. Using $\norm{q}_\infty$
instead would incorrectly claim the upper bound $0.2$.
\paragraph{Why ECDR needs it.}
A margin direction and a perturbation budget become a scalar bound on
possible margin loss. The dual pairing determines the correct constant;
choosing two matching-looking subscripts can be unsound.
\paragraph{Conditions and connections.}
Equality is possible for the complete norm ball. Additional structural
constraints can reduce the attainable loss, so the same inequality may
be conservative. Nonuniform coordinate budgets lead to a weighted
absolute-value sum rather than one common radius.

\concept{F23}{Metrics, balls, and neighborhoods}{F03, F20}
\paragraph{Foundation and definition.}
A metric is a distance function that is nonnegative, vanishes exactly
when its two arguments agree, is symmetric, and satisfies the triangle
inequality. Every norm supplies a metric by taking the norm of a
difference. A ball is the set of points within a specified distance
from one center.
\paragraph{Symbols and meaning.}
Let $x,y,z\in\R^d$, with integer $d\ge1$ counting coordinates.
Write $d_p(x,y)=\norm{x-y}_p$ for distance under norm $L_p$.
For center $x_0\in\R^d$ and radius $\varepsilon>0$, define
\[
 B_p(x_0,\varepsilon)=\{x:d_p(x,x_0)<\varepsilon\},\qquad
 \overline B_p(x_0,\varepsilon)=\{x:d_p(x,x_0)\le\varepsilon\}.
\]
The first is open and the second is closed. A neighborhood of $x_0$
contains some open ball around $x_0$.
\paragraph{Derivation or example.}
For $x_0=(0,0)$ and $\varepsilon=1$, the Euclidean ball is a disk;
the maximum-coordinate ball is a square. The point $(1,0)$ is in both
closed balls but neither open ball. The triangle inequality gives
$d_p(x,z)\le d_p(x,y)+d_p(y,z)$, so every point within radius
$r$ of $y$ lies within radius $r+d_p(y,x_0)$ of anchor $x_0$.
Here $r\ge0$ is another radius, measured in the same coordinate units.
\paragraph{Why ECDR needs it.}
Local certification needs an exact meaning of ``near the anchor.''
Triangle inequalities can transfer distance allowances, but the enlarged
set must still fit inside a region with a valid decision proof.
\paragraph{Conditions and connections.}
Proximity alone does not imply equal class labels. A decision boundary
can pass arbitrarily close to an input. The norm, normalization, and
boundary convention belong to the certificate's contract. F26 formalizes
local change; F28 connects continuity with positive margin.

\concept{F24}{Intervals, Cartesian products, and boxes}{F03, F06}
\paragraph{Foundation and definition.}
A closed real interval includes every scalar between its two endpoints,
including the endpoints. A Cartesian product collects tuples whose
coordinates belong to specified sets. A box is the Cartesian product
of coordinate intervals. It explicitly permits every listed coordinate
combination, whether or not those combinations appeared in a dataset.
\paragraph{Symbols and meaning.}
Let $d\ge1$ be an integer and $\ell,u\in\R^d$ be finite lower and
upper endpoints satisfying $\ell_i\le u_i$ for indices
$i=1,\ldots,d$. The box is
\[
 [\ell,u]=\prod_{i=1}^{d}[\ell_i,u_i]
 =\{x\in\R^d:\ell_i\le x_i\le u_i\text{ for every }i\}.
\]
An equivalent anchor representation is $x=x_0+\delta$ with
$|\delta_i|\le b_i$, where $x_0=(\ell+u)/2$ and
$b=(u-\ell)/2\in\R^d$ has nonnegative entries.
\paragraph{Derivation or example.}
The intervals $[-1,1]$ and $[2,4]$ form a rectangular box in $\R^2$.
For $f(x)=3x_1-2x_2$, minimizing each independent coordinate term
gives $\min f=3(-1)-2(4)=-11$, attained at $(-1,4)$.
The negative coefficient chooses the upper endpoint, illustrating why
order reversal matters in an affine lower bound.
\paragraph{Why ECDR needs it.}
Boxes support simple coordinatewise membership and efficient interval
propagation. They can overapproximate correlated perturbations by
admitting combinations the data seldom produces; verification must
still cover all admitted combinations.
\paragraph{Conditions and connections.}
Zero-width intervals are legitimate and fix coordinates. Positive widths
are needed when a later argument assumes independent variability in
every coordinate. A box is nonempty under the endpoint inequalities;
reversing an endpoint pair makes its product empty.

\concept{F25}{Infimum and supremum versus minimum and maximum}{F06, F07}
\paragraph{Foundation and definition.}
The infimum is the greatest lower bound of a real-valued set, while the
supremum is its least upper bound. A minimum or maximum additionally
must be attained by a member. These definitions separate a best bound
from an actual optimizing input.
\paragraph{Symbols and meaning.}
Let $f:\mathcal S\to\R$ be a scalar function on a nonempty set
$\mathcal S$. Write
$\inf_{x\in\mathcal S}f(x)$ for the infimum of its output set.
A point $x_*\in\mathcal S$ is a minimizer when
$f(x_*)\le f(x)$ for all $x\in\mathcal S$; its value is
$\min_{x\in\mathcal S}f(x)$. The star identifies a candidate
optimizer and does not denote multiplication.
\paragraph{Derivation or example.}
For $f(x)=x$ on $\mathcal S=(0,1)$, the infimum is $0$ but no
minimum exists. Every permitted input has a positive output, yet no
single positive lower bound works throughout the interval. On
$[0,1]$, the minimum exists and equals zero at $x_*=0$.
For a valid enclosure $L\le f(x)$ for every permitted $x$, one has
\[
 L\le\inf_{x\in\mathcal S}f(x).
\]
Here $L\in\R$ is the verifier's lower bound, not necessarily the
exact infimum and not necessarily an attained value.
\paragraph{Why ECDR needs it.}
A positive sound lower bound proves all margins are positive. A loose
nonpositive bound does not identify a violating input. Separating the
bound, true infimum, and optimizer prevents that logical mistake.
\paragraph{Conditions and connections.}
For unbounded-below output sets the infimum may be $-\infty$ in the
extended reals. Empty domains need a separately declared convention and
are excluded here. F27--F28 give sufficient conditions for finite,
attained extrema.

\concept{F26}{Continuity and the epsilon--delta definition}{F23}
\paragraph{Foundation and definition.}
Continuity means that a sufficiently small input change produces any
prescribed small output change. It is a local property of a function,
defined through distances. It does not mean that the function is
constant nearby or that one global rate of change has been proved.
\paragraph{Symbols and meaning.}
Let $f:\mathcal S\subseteq\R^d\to\R^m$ be a function, and let
$x_0\in\mathcal S$. Positive integers $d,m$ count input and output
coordinates. Continuity at $x_0$ means
\[
 \forall\varepsilon>0\;\exists\delta>0\;\forall x\in\mathcal S:
 \norm{x-x_0}_2<\delta
 \Rightarrow\norm{f(x)-f(x_0)}_2<\varepsilon.
\]
The positive scalar $\varepsilon$ is an output tolerance and $\delta$
is an input tolerance, each in its respective coordinate units.
\paragraph{Derivation or example.}
For the scalar function $f(x)=2x+1$, choose
$\delta=\varepsilon/2$. Then
$|f(x)-f(x_0)|=2|x-x_0|<\varepsilon$ whenever
$|x-x_0|<\delta$. If a continuous scalar margin
$g:\mathcal S\to\R$ satisfies $g(x_0)>0$, continuity with output
tolerance $g(x_0)/2$ supplies some neighborhood where
$g(x)>g(x_0)/2>0$. This proves local existence of
a safe neighborhood, without computing a usable radius.
\paragraph{Why ECDR needs it.}
Affine layers and ReLU are continuous; finite compositions therefore
produce continuous logits. A strictly positive winning margin at an
anchor supports a local invariance argument. Quantitative bounds are
needed to turn that existence statement into an executable gate.
\paragraph{Conditions and connections.}
The discrete class decision can be discontinuous at ties even when
logits are continuous. Continuity on a noncompact region does not alone
give a uniform positive margin. F28 adds the relevant compactness and
strict-positivity conditions.

\concept{F27}{Compactness in finite-dimensional Euclidean spaces}{F23, F24, F26}
\paragraph{Foundation and definition.}
A set is compact if every collection of open sets covering it has a
finite subcollection that still covers it. In finite-dimensional
Euclidean space, the Heine--Borel theorem says compactness is equivalent
to being closed and bounded. This equivalence is a theorem specific to
the setting, rather than a universal definition for arbitrary spaces.
\paragraph{Symbols and meaning.}
Let $\mathcal S\subseteq\R^d$, where integer $d\ge1$ counts
coordinates. Boundedness means there is a finite scalar $R\ge0$ with
$\norm{x}_2\le R$ for every $x\in\mathcal S$. Closedness means
that limits of convergent sequences of members remain in $\mathcal S$.
An open cover is a family of open sets whose union contains $\mathcal S$.
\paragraph{Derivation or example.}
A finite-endpoint box $[\ell,u]$ is closed and bounded, hence compact.
For $\ell,u\in\R^d$ with $\ell_i\le u_i$, boundedness follows from
$|x_i|\le\max(|\ell_i|,|u_i|)$ for every coordinate $i$.
The interval $(0,1)$ is bounded but not closed: the sequence
$x_n=1/(n+1)$, for positive integers $n$, lies in $(0,1)$ and converges to zero, outside the
interval. The line $\R$ is closed but unbounded. Neither set is compact.
\paragraph{Why ECDR needs it.}
Bounded closed perturbation domains make worst-case margin optimization
well behaved when margins are continuous. Compactness is the bridge
from infinitely many inputs to existence of an attained worst case;
it does not replace the computation of that worst case.
\paragraph{Conditions and connections.}
The empty set is compact, but useful margin statements should separately
require a nonempty domain. Closed and bounded is insufficient in
general infinite-dimensional normed spaces. ECDR's finite tensor inputs
fit the finite-dimensional hypothesis. F28 states the resulting extrema
and uniform-positivity theorem precisely.

\concept{F28}{Attained extrema and a uniform positive margin}{F01, F25, F26, F27}
\paragraph{Foundation and definition.}
The extreme-value theorem states that a continuous real-valued function
on a nonempty compact set attains a finite minimum and maximum. A useful
consequence is that a continuous function which is strictly positive
everywhere on such a set has one uniform positive lower bound.
\paragraph{Symbols and meaning.}
Let $\mathcal S\subseteq\R^d$ be nonempty and compact. Let
$z_j:\mathcal S\to\R$ be continuous logit functions for class
indices $j=1,\ldots,K$, with integers $d\ge1$ and $K\ge2$.
Fix class $c$ and define each competing margin
$m_{cj}(x)=z_c(x)-z_j(x)$ for $j\ne c$.
\paragraph{Derivation or example.}
Assume $m_{cj}(x)>0$ for every $x\in\mathcal S$ and competitor
$j\ne c$. Each margin attains its minimum at some
$x_j\in\mathcal S$, so $\eta_j=m_{cj}(x_j)>0$. There are only
$K-1$ competitors; therefore
\[
 \eta=\min_{j\ne c}\eta_j>0,
 \qquad m_{cj}(x)\ge\eta
 \quad\text{for every }x\in\mathcal S, j\ne c.
\]
For $m(x)=2-x$ on $[0,1]$, the minimum is $1$. Dropping closedness
can destroy the conclusion: $m(x)=x$ is positive on $(0,1)$ but has
infimum zero.
\paragraph{Why ECDR needs it.}
This conditional theorem explains why strict decision dominance on a
closed bounded region has a positive clearance. A verifier attempts to
compute a sound usable lower bound; the existence theorem alone does
not give that numerical result or an efficient algorithm.
\paragraph{Conditions and connections.}
Tie-consistent decision constancy is weaker than strict dominance: a
chosen class may win ties while its margin equals zero. Neither
empirical agreement nor an unproved continuity assumption can replace
the theorem's hypotheses. A valid finite-precision certificate needs
additional arithmetic-error control.

\concept{F29}{Probability spaces and probability axioms}{F03, F09}
\paragraph{Foundation and definition.}
A probability model consists of a sample space, a collection of
measurable events, and a probability measure. The probability axioms
are nonnegativity, normalization of the whole space to one, and
countable additivity for pairwise disjoint events. They describe a
mathematical model; they do not assert that observed windows are
independent draws from that model.
\paragraph{Symbols and meaning.}
Write $(\Omega,\mathcal F,P)$ for the model: $\Omega$ contains
possible outcomes, $\mathcal F$ is a sigma-algebra of events closed
under complements and countable unions, and $P:\mathcal F\to[0,1]$
assigns probabilities. For pairwise disjoint events
$A_1,A_2,\ldots\in\mathcal F$,
\[
 P(\Omega)=1,\qquad
 P\!\left(\bigcup_{i=1}^{\infty}A_i\right)
 =\sum_{i=1}^{\infty}P(A_i).
\]
The integer $i$ indexes events, not time samples unless defined so.
\paragraph{Derivation or example.}
An event $A$ and its complement $A^c=\Omega\setminus A$ partition
the sample space. Additivity gives $P(A^c)=1-P(A)$.
For a uniformly selected member of a finite dataset of $N\ge1$
windows, each dataset index has probability $1/N$. If $M$ indices pass
a gate, its probability under this particular empirical model is $M/N$.
This construction says nothing by itself about future recordings.
\paragraph{Why ECDR needs it.}
Expected cost and reuse rates refer to a distribution of queries.
Different choices, such as uniform windows or uniform subjects, define
different probability models and can produce different reported rates.
\paragraph{Conditions and connections.}
Distributional coverage is not a universal safety certificate.
Probability one may still allow a nonempty zero-probability exception
set. F30 explains conditioning and independence; F31 defines expected
cost after a probability model has been specified.

\concept{F30}{Conditional probability and independence}{F29}
\paragraph{Foundation and definition.}
Conditioning restricts a probability calculation to an event that has
occurred. For events $A$ and $B$ with $P(B)>0$, define
$P(A\mid B)=P(A\cap B)/P(B)$. Events are independent when
$P(A\cap B)=P(A)P(B)$. Independence is an additional property of
the model; overlapping windows do not acquire it by being placed in
different rows of an array.
\paragraph{Symbols and meaning.}
Let $A$ be the event ``a query lies in an admitted domain'' and $B$ the
event ``the current anchor has a successful valid certificate.''
The operator $P$ is the chosen query-distribution probability, and
$A\cap B$ means both events occur. These events describe distinct
gates and can be strongly dependent.
\paragraph{Derivation or example.}
If $P(B)=0.4$ and $P(A\mid B)=0.6$, then
\[
 P(A\cap B)=P(A\mid B)P(B)=0.24.
\]
Substituting an unconditional estimate $P(A)=0.7$ would give $0.28$,
which is unjustified unless independence or another valid argument is
available. For a finite sample, a conditional rate uses only the
eligible denominator: $60$ accepted queries among $100$ queries with
certified anchors gives $0.6$, even if $250$ total queries were seen.
\paragraph{Why ECDR needs it.}
Certified reuse requires all applicable gates together. Multiplying
independently reported coverage and certification percentages can give
the wrong total reuse rate when their denominators or conditions differ.
\paragraph{Conditions and connections.}
The elementary conditional formula is undefined when $P(B)=0$.
More advanced conditional distributions require additional machinery
and are not supplied by writing a ratio with zero denominator.
Statistical independence differs from orthogonality and zero
correlation. F32--F33 explain empirical weights and covariance.

\concept{F31}{Expectation and the linearity of expected cost}{F09, F29}
\paragraph{Foundation and definition.}
A random variable is a measurable numerical function of an outcome.
Its expectation is a probability-weighted average, formalized as an
integral. For a finite discrete model, this reduces to a weighted sum.
Linearity of expectation does not require independent variables, but
finite ordinary values require suitable integrability.
\paragraph{Symbols and meaning.}
Let $C:\Omega\to\R$ be a random cost measured in microseconds.
For possible costs $c_i$ and probabilities $p_i\ge0$, indexed by
$i=1,\ldots,n$, with integer $n\ge1$ and $\sum_i p_i=1$, define
$\E[C]=\sum_i p_i c_i$. Integrability means
$\E[|C|]<\infty$. The symbol $\E$ refers to the declared probability
model, not an automatic average over an unspecified benchmark.
\paragraph{Derivation or example.}
Let $I$ indicate reuse, so $I=1$ on reuse and $I=0$ otherwise, with
$P(I=1)=r$. Suppose every query pays fixed gate cost $g$ and a miss
pays fixed forward cost $f$, both in microseconds. Then
\[
 C=g+(1-I)f,\qquad \E[C]=g+(1-r)f.
\]
For $g=10$, $f=100$, and $r=0.7$, expected query cost is $40$.
If per-query build cost $B$ is random and integrable, the expectation
becomes $40+\E[B]$. These are hypothetical numbers, not measurements.
\paragraph{Why ECDR needs it.}
Saving forward calls does not establish positive net savings. Gate,
build, failure, and fallback costs must appear in the same expected
cost calculation, with their actual conditional behavior when variable.
\paragraph{Conditions and connections.}
Medians and high quantiles cannot generally replace expectations in
this identity. Unbounded nonintegrable variables require extended-real
care. Correlation between costs does not break finite-sum linearity;
it does affect variance and uncertainty estimates.

\concept{F32}{Empirical means, variances, and weighting}{F09, F31}
\paragraph{Foundation and definition.}
An empirical mean summarizes observed numbers using declared weights.
An empirical population variance measures average squared deviation
under those same weights. These are descriptive definitions. Calling a
variance estimator unbiased requires a sampling model and additional
assumptions, commonly independent identically distributed observations.
\paragraph{Symbols and meaning.}
Let $x_1,\ldots,x_n\in\R$ be measurements, with integer $n\ge1$,
and let $w_i\ge0$ satisfy $\sum_{i=1}^{n}w_i=1$. Define
\[
 \bar x=\sum_{i=1}^{n}w_i x_i,
 \qquad v=\sum_{i=1}^{n}w_i(x_i-\bar x)^2.
\]
The mean $\bar x$ has the measurements' units; variance $v$ has squared
units. Equal weighting uses $w_i=1/n$. Standard deviation is $\sqrt v$.
\paragraph{Derivation or example.}
For equal-weight values $(1,3)$, the mean is $2$ and population
variance is $1$. A recording with one value $0$ and another recording
with nine values $10$ have pooled-sample mean $9$, but an equal-recording
mean of their means is $5$. Both are correct for different weights.
For equal-weight data with $n>1$, the statistic
$s^2=\sum_i(x_i-\bar x)^2/(n-1)$ is a different quantity from $v$.
Under the usual independent identical finite-variance sampling model,
$s^2$ estimates population variance without bias.
\paragraph{Why ECDR needs it.}
Pooling windows weights long recordings more heavily, and overlapping
windows count some raw samples repeatedly. Training normalization and
reported averages must state whether they weight samples, windows,
recordings, or subjects.
\paragraph{Conditions and connections.}
Descriptive means remain defined for dependent windows; ordinary
independent-sample uncertainty formulas may fail. F33 requires a
consistent weighting convention when reporting covariance or correlation.

\concept{F33}{Covariance and Pearson correlation}{F18, F30, F32}
\paragraph{Foundation and definition.}
Covariance measures whether paired variables deviate in the same
direction from their means. Pearson correlation divides covariance by
both standard deviations to obtain a dimensionless measure of linear
association. It is undefined when either variance is zero; assigning
zero in that case is a software policy, not the mathematical definition.
\paragraph{Symbols and meaning.}
Let $(x_i,y_i)$, for $i=1,\ldots,n$, be paired real measurements with
integer $n\ge1$ and normalized nonnegative weights $w_i$. Let
$\bar x,\bar y$ be their weighted means, and let $v_x,v_y$ be their
weighted population variances. Define
\[
 c_{xy}=\sum_i w_i(x_i-\bar x)(y_i-\bar y),
 \qquad \rho=\frac{c_{xy}}{\sqrt{v_xv_y}}
 \quad\text{when }v_xv_y>0.
\]
Covariance has the product of the two measurement units; correlation
$\rho$ has no units.
\paragraph{Derivation or example.}
For equally weighted pairs $(1,2)$ and $(3,6)$,
$\bar x=2$, $\bar y=4$, $v_x=1$, $v_y=4$, and $c_{xy}=2$,
so $\rho=1$. Nonetheless, the vectors $(1,3)$ and $(2,6)$ are
different. Zero covariance need not imply independence: for a uniformly
distributed scalar $X$ on $[-1,1]$, $Y=X^2$ is determined by $X$,
yet symmetry gives zero covariance.
\paragraph{Why ECDR needs it.}
Lag-one raw correlation, lag-64 correlation, within-window shape
correlation, and correlation of window means concern different paired
objects. Their values answer different questions and cannot substitute
for a perturbation norm or decision certificate.
\paragraph{Conditions and connections.}
Specify pair construction, lag, weighting, and treatment of constant or
invalid segments. High correlation can coexist with large absolute
differences because scaling and offsets affect distance differently.

\concept{F34}{Quantiles and empirical distribution functions}{F06, F29, F32}
\paragraph{Foundation and definition.}
A quantile locates a probability threshold in an ordered distribution.
For a finite dataset, the empirical cumulative distribution function
counts how many values lie at or below each threshold. Several
interpolation conventions exist; choosing one is part of a reproducible
analysis and should not be confused with a theorem about future data.
\paragraph{Symbols and meaning.}
For real observations $a_1,\ldots,a_n$ with integer $n\ge1$, define
\[
 F_n(t)=\frac1n\sum_{i=1}^{n}\mathbf1\{a_i\le t\},
 \qquad q_p=\inf\{t\in\R:F_n(t)\ge p\},\quad 0<p\le1.
\]
The scalar $t$ is a threshold; $p$ is a probability level; the indicator
$\mathbf1$ equals one when its condition holds and zero otherwise.
The quantile $q_p$ uses the inverse empirical-CDF convention.
\paragraph{Derivation or example.}
For sorted observations $(1,2,4,8)$, this convention gives
$q_{0.75}=4$. A coordinate budget can use
$a_i=|\delta_{ij}|$, where $\delta_{ij}$ is coordinate $j$ of
observed perturbation $i$. Even if every coordinate has $90\%$
marginal coverage, their joint box coverage need not be $90\%$.
For two coordinates, if their $10\%$ failure events are disjoint,
joint success is $80\%$; if independent, it is $81\%$.
\paragraph{Why ECDR needs it.}
Quantiles provide empirically motivated domain budgets. Actual joint
membership must be measured or bounded separately, and decision safety
still requires a whole-domain proof. A large observed coverage value
does not supply that proof.
\paragraph{Conditions and connections.}
Budgets fitted on training data need a stated fitting population and
evaluation split. Dependence, finite-sample uncertainty, distribution
shift, and quantile interpolation affect interpretation. Quantiles of
raw and normalized coordinates are not interchangeable.



